
\subsection*{1. The two-dimensional determinant}
For columns \(v=(a,c)^\top\) and \(w=(b,d)^\top\), set
\[
 \det(v,w)=\begin{vmatrix}a&b\\c&d\end{vmatrix}=ad-bc.
\]
This expression is linear in each column separately, vanishes when the
columns coincide, and takes the value \(1\) on \((e_1,e_2)\).

\begin{proposition}[Characterization in dimension two]
If \(f:\R^2\times\R^2\to\R\) is linear in each argument,
\(f(v,v)=0\), and \(f(e_1,e_2)=1\), then \(f(v,w)=\det(v,w)\).
\end{proposition}
\begin{proof}
Expanding \(0=f(v+w,v+w)\) gives \(f(v,w)=-f(w,v)\).
Consequently
\[
 f(ae_1+ce_2,be_1+de_2)
 =abf(e_1,e_1)+adf(e_1,e_2)+cbf(e_2,e_1)+cdf(e_2,e_2)
 =ad-bc .
\]
\end{proof}
Thus exchanging columns reverses the sign; adding a multiple of one column
to the other leaves the determinant unchanged. These are consequences of
linearity and vanishing on repeated columns, not special computational tricks.

\begin{example}
For \(A=\begin{pmatrix}2&1\\1&4\end{pmatrix}\), \(\det A=8-1=7\).
Replacing its second column by the second column minus twice the first gives
\(\begin{pmatrix}2&-3\\1&2\end{pmatrix}\), still with determinant \(4+3=7\).
The vectors \((2,1)^\top,(1,4)^\top\) are independent: taking determinants of
\(x(2,1)^\top+y(1,4)^\top=0\) against each column gives \(7x=7y=0\).
\end{example}

\subsection*{2. Cofactors and determinant operations}
For a square matrix \(A=(a_{ij})\), let \(A_{ij}\) be the matrix obtained by
deleting row \(i\) and column \(j\). Its \emph{minor} is \(\det A_{ij}\),
and its \emph{cofactor} is \(c_{ij}=(-1)^{i+j}\det A_{ij}\).
Set \(\det[\,]=1\) for the empty matrix and \(\det[a]=a\).
The recursive first-row definition is
\[
 \det A=\sum_{j=1}^n a_{1j}(-1)^{1+j}\det A_{1j}.
\]
For \(n=3\) it gives
\[
 \det A=a_{11}a_{22}a_{33}+a_{12}a_{23}a_{31}+a_{13}a_{21}a_{32}
       -a_{13}a_{22}a_{31}-a_{12}a_{21}a_{33}-a_{11}a_{23}a_{32}.
\]

\begin{theorem}[Determinant laws]\label{thm:07:laws}
The determinant is the unique function of \(n\) columns that is linear in
each column, vanishes on repeated columns, and equals \(1\) on \(I_n\).
Moreover,
\[
 \det A^\top=\det A,\qquad
 \det A=\sum_j a_{ij}c_{ij}=\sum_i a_{ij}c_{ij}
\]
for each fixed row \(i\) and each fixed column \(j\), respectively.
\end{theorem}
The complete construction and proof are in the reading supplement below.
The following consequences will be used in every calculation.

\begin{proposition}[Row and column operations]
Exchanging two rows changes the sign; multiplying one row by \(c\) multiplies
the determinant by \(c\); adding a multiple of another row leaves it
unchanged. A zero row or two equal rows give determinant zero.
The same statements hold for columns. In particular,
\[
 \det(cA)=c^n\det A.
\]
For a triangular matrix, \(\det A=\prod_{i=1}^n a_{ii}\).
\end{proposition}
\begin{proof}
For columns, expand by linearity. The sign rule follows by expanding the
determinant with two selected columns both equal to \(u+v\); the two repeated
terms vanish. A row statement follows by transposition. Scaling the whole
matrix scales each of its \(n\) rows. For an upper triangular matrix, expand
in its first column and repeat. For a lower triangular matrix, transpose.
\end{proof}

\begin{example}[A cofactor expansion and an elimination calculation]
For
\[
 A=\begin{pmatrix}3&0&1\\1&2&5\\-1&4&2\end{pmatrix},
\]
expansion in column \(2\) gives
\[
 \det A=2\begin{vmatrix}3&1\\-1&2\end{vmatrix}
       -4\begin{vmatrix}3&1\\1&5\end{vmatrix}
       =2(7)-4(14)=-42.
\]
Alternatively \(R_3\leftarrow R_3-2R_2\) gives
\(\det A=2\begin{vmatrix}3&1\\-3&-8\end{vmatrix}=2(-21)=-42\).
For
\[
 B=\begin{pmatrix}1&3&1&1\\2&1&5&2\\1&-1&2&3\\4&1&-3&7\end{pmatrix},
\]
use \(R_1\leftarrow R_1+3R_3\), \(R_2\leftarrow R_2+R_3\),
\(R_4\leftarrow R_4+R_3\). Expansion in column \(2\) then gives
\[
 \det B=\begin{vmatrix}4&7&10\\3&7&5\\5&-1&10\end{vmatrix}
 =4(75)-7(5)+10(-38)=-115.
\]
The coefficient is \(+1\): the entry \(-1\) in position \((3,2)\) and its
cofactor sign \((-1)^5\) cancel.
\end{example}

\subsection*{3. Dependence and rank}
\begin{theorem}[Nonsingularity criterion]\label{thm:07:nonsingular}
For a square matrix \(A\), the following are equivalent:
\(\det A\ne0\); its columns are independent; \(A\) is invertible;
\(Ax=b\) has exactly one solution for every \(b\).
\end{theorem}
\begin{proof}
The last three equivalences are the square-system theorem.
If a column is a combination of the others, multilinearity expresses
\(\det A\) as a sum of determinants with repeated columns, hence zero.
Conversely, if the columns are independent, elimination gives an upper
triangular matrix with all diagonal entries nonzero. Each elementary
operation changes the determinant by a nonzero factor, and the final
determinant is the product of its nonzero diagonal entries. Thus
\(\det A\ne0\).
\end{proof}

For a rectangular matrix, a minor of order \(r\) is the determinant of a
submatrix obtained by selecting \(r\) rows and \(r\) columns.

\begin{proposition}[Rank from minors]
The rank of \(A\) is the largest order of a nonzero minor. The zero matrix
has rank zero.
\end{proposition}
\begin{proof}
A nonzero minor of order \(s\) makes its \(s\) selected columns independent,
even after the discarded coordinates are restored; thus \(s\leqslant
\operatorname{rank}A\). If the rank is \(r>0\), select \(r\) independent
columns and call the resulting matrix \(C\). Row rank equals column rank,
so \(C\) has \(r\) independent rows. Selecting them gives a nonsingular
\(r\times r\) matrix, hence a nonzero minor of \(A\).
\end{proof}

\begin{example}[A zero minor does not determine the rank]
The first three columns of
\[
 A=\begin{pmatrix}3&1&2&5\\1&2&-1&2\\4&3&1&1\end{pmatrix}
\]
have determinant zero, since the last row in that submatrix is the sum of
the first two. But the last three columns have determinant
\[
 \begin{vmatrix}1&2&5\\2&-1&2\\3&1&1\end{vmatrix}
 =-3+8+25=30.
\]
Thus \(\operatorname{rank}A=3\).
If its bottom-right entry is changed to \(7\), the entire third row is the
sum of the first two. The rank is then \(2\), since
\(\begin{vmatrix}3&1\\1&2\end{vmatrix}=5\ne0\).
\end{example}

\subsection*{4. Cramer's rule and the inverse}
\begin{theorem}[Cramer's rule]
Suppose \(\det A\ne0\). Let \(A[j:b]\) denote \(A\) with its \(j\)-th
column replaced by \(b\). The solution of \(Ax=b\) satisfies
\[
 x_j=\frac{\det A[j:b]}{\det A}.
\]
\end{theorem}
\begin{proof}
Write \(b=\sum_k x_k a_k\), where \(a_k\) is column \(k\) of \(A\).
Expand \(\det A[j:b]\) in column \(j\). Every term except \(k=j\)
has a repeated column and vanishes, leaving \(x_j\det A\).
\end{proof}

\begin{example}
For
\[
 \begin{pmatrix}3&2&4\\2&-1&1\\1&2&3\end{pmatrix}
 \begin{pmatrix}x\\y\\z\end{pmatrix}
 =\begin{pmatrix}1\\0\\1\end{pmatrix},
\]
the denominator is
\(3(-5)-2(5)+4(5)=-5\). The three numerator determinants are
\[
 \begin{vmatrix}1&2&4\\0&-1&1\\1&2&3\end{vmatrix}=1,\quad
 \begin{vmatrix}3&1&4\\2&0&1\\1&1&3\end{vmatrix}=0,\quad
 \begin{vmatrix}3&2&1\\2&-1&0\\1&2&1\end{vmatrix}=-2.
\]
Therefore \((x,y,z)=(-1/5,0,2/5)\). Substitution gives
\((-3/5+8/5,-2/5+2/5,-1/5+6/5)=(1,0,1)\).
\end{example}

Let \(C=(c_{ij})\) be the cofactor matrix and define
\(\operatorname{adj}A=C^\top\), the \emph{adjugate}.

\begin{proposition}[Adjugate identities]\label{prop:07:adj}
For every square \(A\), including singular matrices,
\[
 A\operatorname{adj}A=\operatorname{adj}A\,A=(\det A)I.
\]
If \(\det A\ne0\), then \(A^{-1}=(\det A)^{-1}\operatorname{adj}A\).
\end{proposition}
\begin{proof}
Entry \((i,k)\) of \(AC^\top\) is \(\sum_j a_{ij}c_{kj}\).
For \(i=k\), this is the cofactor expansion of \(\det A\).
For \(i\ne k\), it is the expansion in row \(k\) of the matrix obtained by
replacing row \(k\) by row \(i\): the cofactors in that row do not change,
and the repeated rows make the determinant zero. The column version
proves \(C^\top A=(\det A)I\). Division gives the inverse formula when
\(\det A\ne0\).
\end{proof}

\begin{example}[Compute every column of the inverse]
For
\[
 A=\begin{pmatrix}3&1&-2\\-1&1&2\\1&-2&1\end{pmatrix},
\qquad
 C=\begin{pmatrix}5&3&1\\3&5&7\\4&-4&4\end{pmatrix},
\]
the first row gives \(\det A=3(5)+1(3)-2(1)=16\).
For instance, \(c_{23}=-\det\begin{pmatrix}3&1\\1&-2\end{pmatrix}=7\).
Thus
\[
 A^{-1}=\frac1{16}\begin{pmatrix}5&3&4\\3&5&-4\\1&7&4\end{pmatrix}.
\]
Multiplication gives
\[
 A\begin{pmatrix}5&3&4\\3&5&-4\\1&7&4\end{pmatrix}
 =\begin{pmatrix}16&0&0\\0&16&0\\0&0&16\end{pmatrix}.
\]
The transpose of the cofactor matrix is essential.
\end{example}

\subsection*{5. Products and block elimination}
\begin{theorem}[Multiplicativity]\label{thm:07:product}
For square matrices \(A,B\) of the same order,
\[
 \det(AB)=\det A\,\det B.
\]
Consequently \(\det A^{-1}=1/\det A\) for invertible \(A\), and an
orthogonal matrix has determinant \(1\) or \(-1\).
\end{theorem}
\begin{proof}
Fix \(A\) and put \(f(b_1,\ldots,b_n)=\det(Ab_1,\ldots,Ab_n)\).
This is alternating and multilinear. Expanding the columns in standard
coordinates, as in the uniqueness proof in the supplement, shows that
\(f(b_1,\ldots,b_n)=f(e_1,\ldots,e_n)\det B\).
Here \(f(e_1,\ldots,e_n)=\det A\), including when it is zero.
The inverse formula follows from \(\det I=1\). If \(Q^\top Q=I\),
then \(1=\det Q^\top\det Q=(\det Q)^2\).
\end{proof}

\begin{proposition}[Block triangular and Schur complement formulas]
For conformable blocks with \(A\) of order \(m\) and \(D\) of order \(n\),
\[
 \det\begin{pmatrix}A&B\\0&D\end{pmatrix}=\det A\det D.
\]
If \(A\) is invertible, then
\[
 \det\begin{pmatrix}A&B\\C&D\end{pmatrix}
 =\det A\,\det(D-CA^{-1}B).
\]
Neither \(D\) nor the whole block matrix is assumed invertible.
\end{proposition}
\begin{proof}
In the permutation formula, a nonzero term of the block triangular
determinant must assign the last \(n\) rows to the last \(n\) columns;
the remaining rows go to the first \(m\) columns. The sum factors into the
two determinant sums. Transposition gives the lower triangular formula.
For the second assertion, multiply on the left by
\(\begin{pmatrix}I&0\\-CA^{-1}&I\end{pmatrix}\), whose determinant is \(1\).
The result is
\(\begin{pmatrix}A&B\\0&D-CA^{-1}B\end{pmatrix}\).
\end{proof}

\subsection*{6. Area, volume, and interpolation}
\begin{proposition}[Geometric meaning]
If \(v,w\in\R^2\), the area of their parallelogram is
\(|\det(v,w)|\). If \(v_1,\ldots,v_n\in\R^n\), the volume of their
parallelotope is \(|\det(v_1,\ldots,v_n)|\).
Consequently a linear map with matrix \(T\) multiplies such volumes by
\(|\det T|\).
\end{proposition}
\begin{proof}
In the plane, the case \(v=0\) is immediate. Otherwise write
\(v=(a,c)^\top\), \(w=(b,d)^\top\), and take the unit normal
\(u=(-c,a)^\top/\norm v\). The signed height of \(w\) above the directed
line through \(v\) is \(u^\top w=(ad-bc)/\norm v\).
Base times absolute height is therefore \(|ad-bc|\).

For independent vectors in \(\R^n\), Gram--Schmidt gives \(V=QR\), with
\(Q\) orthogonal and \(R\) upper triangular with positive diagonal.
The successive perpendicular heights are \(r_{11},\ldots,r_{nn}\);
base times height recursively gives volume \(\prod_i r_{ii}\).
Since \(|\det Q|=1\), this equals \(|\det V|\).
Dependent vectors lie in a proper subspace and have zero \(n\)-dimensional
volume, as does their determinant. Finally
\(|\det(TV)|=|\det T|\,|\det V|\).
\end{proof}
Here the geometric base-times-height rule, and the fact that a region in a
proper subspace has zero ambient volume, are the usual Euclidean volume
facts. A derivation from signed planar area is also supplied below.

\begin{example}
The parallelogram with vertices \(p=(1,1)\), \(q=(2,-1)\),
\(r=(4,6)\), and \(q+r-p\) has adjacent sides
\(q-p=(1,-2)\), \(r-p=(3,5)\); its area is \(|5+6|=11\).
The map \(T=\begin{pmatrix}2&1\\1&4\end{pmatrix}\) takes these sides to
\((0,-7)\), \((11,23)\), whose determinant is \(77=7\cdot11\).
Translation of the initial point does not change the side vectors.
\end{example}

\begin{example}[Vandermonde and quadratic interpolation]
For distinct \(x_0,x_1,x_2\),
\[
 \det\begin{pmatrix}1&x_0&x_0^2\\1&x_1&x_1^2\\1&x_2&x_2^2\end{pmatrix}
 =(x_1-x_0)(x_2-x_0)(x_2-x_1).
\]
Subtract the first row from each of the other rows, factor
\(x_1-x_0\) and \(x_2-x_0\), and expand in column \(1\).
The remaining determinant is
\(\det\begin{pmatrix}1&x_1+x_0\\1&x_2+x_0\end{pmatrix}=x_2-x_1\).
Thus values \(y_0,y_1,y_2\) at three distinct points determine exactly one
polynomial of degree at most two. For nodes \(-1,0,1\) and values \(2,1,4\),
the equations are \(a_0-a_1+a_2=2\), \(a_0=1\),
\(a_0+a_1+a_2=4\), giving \(p(t)=1+t+2t^2\).
The general determinant and interpolation construction are proved in the
starred exercises.
\end{example}

\subsection*{Reading supplement: construction of the determinant}
A permutation \(\sigma\) of \(\{1,\ldots,n\}\) has sign
\(\operatorname{sgn}\sigma=(-1)^{N(\sigma)}\), where \(N(\sigma)\)
counts pairs \(i<j\) with \(\sigma(i)>\sigma(j)\).
Interchanging two positions reverses this sign: for positions \(r<s\),
the pair \((r,s)\) reverses, while the changes from intermediate positions
occur in pairs. Sorting a permutation by adjacent exchanges proves that
a product of \(k\) exchanges has sign \((-1)^k\); in particular sign
multiplies under composition and
\(\operatorname{sgn}(\sigma^{-1})=\operatorname{sgn}\sigma\).

Define
\[
 D(A)=\sum_{\sigma}\operatorname{sgn}(\sigma)
           \prod_{j=1}^n a_{\sigma(j),j}.
\]
Each factor contains one entry from each column, so \(D\) is multilinear.
Exchanging two columns permutes the summands by an exchange of two
positions and reverses the sign. Thus repeated columns give \(D=-D=0\).
Only the identity permutation contributes to \(D(I)\), giving \(1\).

Conversely, let \(f\) be any alternating multilinear function.
Expanding \(a_j=\sum_i a_{ij}e_i\) gives
\[
 f(a_1,\ldots,a_n)
 =\sum_{i_1,\ldots,i_n}a_{i_1,1}\cdots a_{i_n,n}
                      f(e_{i_1},\ldots,e_{i_n}).
\]
Repeated indices give zero. All remaining index lists are permutations,
and exchanging entries shows that
\(f(e_{\sigma(1)},\ldots,e_{\sigma(n)})
 =\operatorname{sgn}(\sigma)f(e_1,\ldots,e_n)\).
Hence \(f(A)=f(I)D(A)\). This proves uniqueness.

Replacing \(\sigma\) by \(\sigma^{-1}\) in the sum proves
\(D(A^\top)=D(A)\). To expand in entry \(a_{ij}\), group permutations
having \(\sigma(j)=i\). Moving row \(i\) and column \(j\) to the first
positions takes \(i-1+j-1\) exchanges; the resulting sign is
\((-1)^{i+j}\). The remaining permutation sum is \(D(A_{ij})\).
Grouping all choices of \(j\) in fixed row \(i\) gives
\(D(A)=\sum_j(-1)^{i+j}a_{ij}D(A_{ij})\); transpose gives column
expansion. In particular \(D\) satisfies the recursive definition, proving
all claims of Theorem~\ref{thm:07:laws}.

\subsection*{Reading supplement: why signed area is alternating and linear}
Give a parallelogram the sign of the orientation of its ordered side
vectors, and call its signed area \(F(v,w)\). A degenerate parallelogram
has area zero; the standard square has signed area \(1\).
For a positive integer \(k\), the parallelogram on \(kv,w\) is a union
of \(k\) translates of the parallelogram on \(v,w\), with intersections
only along edges. Additivity and translation invariance of area give
\(F(kv,w)=kF(v,w)\). Applying this to \(v/k\) gives rational positive
scaling. For real \(c>0\), approximate \(c\) above and below by positive
rationals. The corresponding parallelograms are nested, so monotonicity
squeezes their unsigned areas to \(c\) times the original area.
Negating one side reverses orientation and preserves unsigned area.
Thus \(F(cv,w)=cF(v,w)\) for every real \(c\), including \(c=0\).
Exchanging the side vectors reverses orientation and preserves area:
\(F(w,v)=-F(v,w)\).

Replacing \(v\) by \(v+tw\) is a shear parallel to the base \(w\).
The base and signed perpendicular height are unchanged, so
\(F(v+tw,w)=F(v,w)\). This also follows by cutting off one triangular
piece and translating it to the other end of the base.

It remains to prove additivity, without assuming it. If \(w=0\), it is
immediate. Choose \(u\) independent of \(w\). Write
\(v=au+bw\) and \(z=cu+dw\). Scaling and shear invariance give
\(F(v,w)=aF(u,w)\), also when \(a=0\); likewise
\(F(z,w)=cF(u,w)\) and \(F(v+z,w)=(a+c)F(u,w)\).
Hence \(F(v+z,w)=F(v,w)+F(z,w)\). Alternation gives linearity in the
second variable. The two-dimensional characterization now yields
\(F(v,w)=\det(v,w)\).

\clearpage
\section*{Exercises and solutions}
A star marks an advanced exercise. In block matrices, all products are
assumed conformable; each problem states any additional size or
invertibility assumptions.
\subsection*{Calculation, dependence, and geometry}

\begin{problem}[Cofactor calculations]\label{ex:07:three}
Compute the determinants of
\[
 \begin{pmatrix}2&1&2\\0&3&-1\\4&1&1\end{pmatrix},\quad
 \begin{pmatrix}2&4&3\\-1&3&0\\0&2&1\end{pmatrix},\quad
 \begin{pmatrix}-1&5&3\\4&0&0\\2&7&8\end{pmatrix},\quad
 \begin{pmatrix}3&1&2\\4&5&1\\-1&2&-3\end{pmatrix}.
\]
\end{problem}
\begin{solution}
Expanding in row \(1\) for the first, second, and fourth, and in row \(2\)
for the third, gives respectively
\[
 2(4)-1(4)+2(-12)=-20,\qquad 2(3)-4(-1)+3(-2)=4,
\]
\[
 -4(40-21)=-76,\qquad 3(-17)-1(-11)+2(13)=-14.
\]
\end{solution}

\begin{problem}[Use the structure]\label{ex:07:four}
Compute the determinants of
\[
 A=\begin{pmatrix}1&1&-2&4\\0&1&1&3\\2&-1&1&0\\3&1&2&5\end{pmatrix},
 \quad B=\begin{pmatrix}-1&1&2&0\\0&3&2&1\\0&4&1&2\\3&1&5&7\end{pmatrix},
\]
\[
 C=\begin{pmatrix}3&1&1\\2&5&5\\8&7&7\end{pmatrix},
 \quad D=\begin{pmatrix}4&-1&1\\2&0&0\\1&5&7\end{pmatrix}.
\]
\end{problem}
\begin{solution}
For \(A\), subtract \(2R_1\) from \(R_3\) and \(3R_1\) from \(R_4\).
Expansion in column \(1\), followed by two row additions, gives
\[
 \det A=\begin{vmatrix}1&1&3\\-3&5&-8\\-2&8&-7\end{vmatrix}
 =\begin{vmatrix}1&1&3\\0&8&1\\0&10&-1\end{vmatrix}=-18.
\]
For \(B\), use \(R_4\leftarrow R_4+3R_1\):
\[
 \det B=-\begin{vmatrix}3&2&1\\4&1&2\\4&11&7\end{vmatrix}
 =-\{3(-15)-2(20)+40\}=45.
\]
The last two columns of \(C\) coincide, so \(\det C=0\).
Expanding \(D\) in row \(2\) gives
\(-2\det\begin{pmatrix}-1&1\\5&7\end{pmatrix}=24\).
\end{solution}

\begin{problem}[A parameter and four operations]\label{ex:07:parameter}
Let \(A_\alpha=\begin{pmatrix}1&\alpha&0\\4&5&3\\1&0&2\end{pmatrix}\).
\begin{enumerate}
\item Find when its determinant is zero and explain the column dependence.
\item Compute the determinant after exchanging rows \(2\) and \(3\).
\item Compute it after doubling column \(1\) of the original matrix.
\item Compute it after subtracting four times row \(1\) from row \(2\) of the original matrix.
\end{enumerate}
\end{problem}
\begin{solution}
Expansion in row \(1\) gives \(10-5\alpha\), so the singular value is
\(\alpha=2\). At this value, \(2a_1-a_2-a_3=0\); otherwise the columns
are independent by the nonsingularity criterion.
The three further determinants are respectively
\(-10+5\alpha\), \(20-10\alpha\), and \(10-5\alpha\).
Each operation in (b)--(d) is applied separately to the original matrix.
\end{solution}

\begin{problem}[A tridiagonal determinant]\label{ex:07:tridiagonal}
Let \(T_n\) have \(2\) on its diagonal, \(-1\) immediately above and below
the diagonal, and zero elsewhere. Prove that \(\det T_n=n+1\).
\end{problem}
\begin{solution}
Put \(d_0=1\), \(d_1=2\). Expansion in the first row gives
\(d_n=2d_{n-1}-d_{n-2}\): the second minor has first column
\((-1,0,\ldots,0)^\top\), so its determinant is \(-d_{n-2}\).
If \(d_{n-1}=n\) and \(d_{n-2}=n-1\), then
\(d_n=2n-(n-1)=n+1\). The initial cases complete the induction.
\end{solution}

\begin{problem}[*Vandermonde determinant]\label{ex:07:vandermonde}
For distinct real numbers \(a_0,\ldots,a_n\), prove
\[
 \det\begin{pmatrix}
 1&1&\cdots&1\\a_0&a_1&\cdots&a_n\\
 \vdots&\vdots&&\vdots\\a_0^n&a_1^n&\cdots&a_n^n
 \end{pmatrix}
 =\prod_{0\leqslant j<i\leqslant n}(a_i-a_j).
\]
\end{problem}
\begin{solution}
For \(n=0\) the empty product is \(1\). Regard the determinant as a
polynomial \(p(t)\) in its last node \(t=a_n\). Its degree is at most \(n\).
Its roots include \(a_0,\ldots,a_{n-1}\), because a repeated node gives
repeated columns. Its coefficient of \(t^n\), obtained by expanding in
the last column, is the Vandermonde determinant on the first \(n\) nodes.
The factor theorem therefore gives
\(p(t)=\det V_{n-1}\prod_{j=0}^{n-1}(t-a_j)\).
Induction proves the formula. The formula also remains valid when nodes
coincide: both sides then vanish.
\end{solution}

\begin{problem}[Inverse, orthogonality, and orientation]\label{ex:07:orientation}
\begin{enumerate}
\item Show that \(\det A^{-1}=1/\det A\) for invertible \(A\), and that
\(\det Q=\pm1\) for a real orthogonal matrix \(Q\).
\item Find the determinants of
\[
 R=\begin{pmatrix}\cos\theta&-\sin\theta\\\sin\theta&\cos\theta\end{pmatrix},
 \qquad
 H=\begin{pmatrix}\cos\theta&-\sin\theta\\-\sin\theta&-\cos\theta\end{pmatrix},
\]
and explain the difference.
\end{enumerate}
\end{problem}
\begin{solution}
Taking determinants of \(AA^{-1}=I\) proves the first identity.
Taking determinants of \(Q^\top Q=I\) gives \((\det Q)^2=1\).
Direct calculation gives \(\det R=1\) and \(\det H=-1\).
Both preserve lengths and areas; \(R\) preserves orientation while \(H\)
reverses it. Indeed \(H=\operatorname{diag}(1,-1)R\), a rotation followed
by reflection in the first coordinate axis.
\end{solution}

\begin{problem}[*Further adjugate identities]\label{ex:07:adjugate}
For square \(A,B\) of order \(n\geqslant1\), prove, without assuming
invertibility,
\[
 \det(\operatorname{adj}A)=(\det A)^{n-1},\qquad
 \operatorname{adj}(AB)=(\operatorname{adj}B)(\operatorname{adj}A).
\]
For \(n=1\), interpret the first right-hand side as \(1\).
\end{problem}
\begin{solution}
If \(A\) is invertible, \(\operatorname{adj}A=(\det A)A^{-1}\), so
\(\det(\operatorname{adj}A)=(\det A)^n/\det A\).
If both matrices are invertible, the same formula and the order of the
product inverse give the second identity.

To include singular matrices, observe that every cofactor is a polynomial
in the entries. A nonzero real polynomial has only finitely many roots;
this follows inductively from the factor theorem. Thus \(A+tI\) is
invertible for all but finitely many \(t\), since its determinant is a
monic polynomial in \(t\). The first identity, applied to \(A+tI\),
is an equality of polynomials in \(t\) for infinitely many \(t\);
their difference is identically zero. Evaluate at \(t=0\).
For the second identity use \(A+tI,B+tI\): outside a finite set both
are invertible. Both sides, including the adjugate of their product,
are polynomial matrices in \(t\); the same argument proves the identity
at zero. For \(n=1\), the adjugate is \([1]\), agreeing with both formulas.
\end{solution}

\begin{problem}[Compute the adjugate inverse]\label{ex:07:inverse}
Find the inverse of \(A=\begin{pmatrix}1&1&0\\4&5&3\\1&0&2\end{pmatrix}\).
\end{problem}
\begin{solution}
Its determinant is \(5\). Computing all nine signed minors gives
\[
 C=\begin{pmatrix}10&-5&-5\\-2&2&1\\3&-3&1\end{pmatrix}.
\]
Hence
\[
 A^{-1}=\frac15\begin{pmatrix}10&-2&3\\-5&2&-3\\-5&1&1\end{pmatrix}.
\]
For a direct check, multiplication by \(A\) gives diagonal entries \(5\)
and off-diagonal entries zero before division by \(5\).
\end{solution}

\begin{problem}[Determinants of elementary matrices]\label{ex:07:elementary}
For order \(n\), find the determinants of the elementary matrices that
exchange distinct rows \(i,j\), scale row \(i\) by \(\lambda\ne0\), or
add \(\lambda\) times row \(j\) to distinct row \(i\).
Explain why the answers do not depend on \(n\).
\end{problem}
\begin{solution}
Apply each operation to \(I_n\), whose determinant is \(1\).
The answers are \(-1,\lambda,1\), respectively, by the row-operation
laws. Each operation concerns the stated rows only; the other unit
rows contribute no further factor. Every elementary matrix therefore
has nonzero determinant.
\end{solution}

\begin{problem}[Two systems by Cramer's rule]\label{ex:07:cramer-pair}
Solve
\[
 \begin{pmatrix}2&1&-3\\5&2&-6\\3&-1&-4\end{pmatrix}x
 =\begin{pmatrix}1\\5\\7\end{pmatrix},\qquad
 \begin{pmatrix}2&1&-2\\3&2&2\\5&4&3\end{pmatrix}y
 =\begin{pmatrix}10\\1\\4\end{pmatrix}
\]
using Cramer's rule.
\end{problem}
\begin{solution}
For the first matrix, expansion in row \(1\) gives
\(2(-14)-1(-2)-3(-11)=7\).
Replacing columns \(1,2,3\) by its right-hand side gives determinants
\(21,-14,7\), respectively. For example the first numerator is
\[
 \begin{vmatrix}1&1&-3\\5&2&-6\\7&-1&-4\end{vmatrix}
 =-14-22+57=21.
\]
Thus \(x=(3,-2,1)^\top\).
For the second matrix, the denominator is
\(2(-2)-1(-1)-2(2)=-7\);
the replacement determinants are \(-7,-14,21\), so \(y=(1,2,-3)^\top\).
These give right-hand sides
\((6-2-3,15-4-6,9+2-4)=(1,5,7)\) and
\((2+2+6,3+4-6,5+8-9)=(10,1,4)\), respectively.
\end{solution}

\begin{problem}[*Interpolation at distinct nodes]\label{ex:07:interpolation}
Given distinct real \(x_0,\ldots,x_n\) and arbitrary real \(y_0,\ldots,y_n\),
find the coefficients of the unique polynomial \(p\) of degree at most
\(n\) satisfying \(p(x_i)=y_i\).
\end{problem}
\begin{solution}
Let \(V=(x_i^j)_{0\leqslant j,i\leqslant n}\), with powers indexing rows.
If \(p(t)=\sum_{j=0}^n a_jt^j\), its coefficients satisfy \(V^\top a=y\).
The Vandermonde determinant is nonzero, so the coefficients exist uniquely:
\[
 a_j=\frac{\det V[\text{row }j+1:y^\top]}{\det V}.
\]
A more explicit form is
\[
 p(t)=\sum_{i=0}^n y_i
       \prod_{\substack{0\leqslant k\leqslant n\\k\ne i}}
                  \frac{t-x_k}{x_i-x_k}.
\]
Each product is \(1\) at \(x_i\) and \(0\) at the other nodes, proving the
interpolation conditions directly; expansion gives every \(a_j\).
If two polynomials worked, their difference, of degree at most \(n\),
would have \(n+1\) distinct roots and hence be zero. The degree can be
less than \(n\), for example when every \(y_i\) is zero.
\end{solution}

\begin{problem}[Rank certificates]\label{ex:07:rank}
Find the ranks of the following matrices, giving both a nonzero minor
and an upper bound when needed:
\[
 A=\begin{pmatrix}2&3&5&1\\1&-1&2&1\end{pmatrix},\quad
 B=\begin{pmatrix}3&5&1&4\\2&-1&1&1\\5&4&2&5\end{pmatrix},
 \quad C=\begin{pmatrix}3&5&1&4\\2&-1&1&1\\7&1&2&5\end{pmatrix},
\]
\[
 D=\begin{pmatrix}-1&1&6&5\\1&1&2&3\\-1&2&5&4\\2&1&0&1\end{pmatrix},
 \quad E=\begin{pmatrix}2&1&6&6\\3&1&1&-1\\5&2&7&5\\8&3&8&4\end{pmatrix}.
\]
\end{problem}
\begin{solution}
For \(A\), the first two columns give determinant \(-5\), so the rank is
\(2\). For \(B\), row \(3=R_1+R_2\), while the top-left minor is
\(-13\); its rank is \(2\). The first three columns of \(C\) have
determinant \(15\), so its rank is \(3\).
For \(D\), add \(R_1\) to \(R_2\), subtract \(R_1\) from \(R_3\), and
add \(2R_1\) to \(R_4\). Expansion in the first column gives
\[
 \det D=-\begin{vmatrix}2&8&8\\1&-1&-1\\3&12&11\end{vmatrix}
 =-\{2(1)-8(14)+8(15)\}=-10,
\]
so its rank is \(4\).
For \(E\), \(R_3=R_1+R_2\) and \(R_4=R_1+2R_2\);
the top-left minor is \(-1\). Its rank is \(2\).
\end{solution}

\begin{problem}[Exponential functions and a determinant]\label{ex:07:exponentials}
Let \(\alpha_1,\ldots,\alpha_n\) be distinct nonzero real numbers.
Prove that \(e^{\alpha_1t},\ldots,e^{\alpha_nt}\), considered as functions
on \(\R\), are linearly independent.
\end{problem}
\begin{solution}
If \(\sum_i c_i e^{\alpha_it}=0\) for every \(t\), differentiate
\(j=0,\ldots,n-1\) times and evaluate at \(0\). This gives
\(\sum_i\alpha_i^j c_i=0\). The coefficient matrix is Vandermonde,
with determinant \(\prod_{i<k}(\alpha_k-\alpha_i)\ne0\).
Therefore every \(c_i=0\). The argument in fact allows one \(\alpha_i=0\).
\end{solution}

\begin{problem}[Three and four unknowns]\label{ex:07:cramer-lang}
Use determinants to solve
\[
 \begin{cases}3x+y-z=0,\\x+y+z=0,\\y-z=1,\end{cases}
 \qquad
 \begin{cases}
 4x+y+z+w=1,\\x-y+2z-3w=0,\\2x+y+3z+5w=0,\\x+y-z-w=2.
 \end{cases}
\]
\end{problem}
\begin{solution}
The first coefficient determinant is \(-6\), and the replacement
determinants are \(2,-4,2\). Thus
\((x,y,z)=(-1/3,2/3,-1/3)\).
For the second, the coefficient determinant is \(48\); replacement
determinants in column order are \(-10,97,16,-25\). Hence
\[
 (x,y,z,w)=\frac1{48}(-10,97,16,-25).
\]
These determinants can be evaluated by row additions before expansion.
For verification, the four numerator combinations are
\(-40+97+16-25=48\),
\(-10-97+32+75=0\),
\(-20+97+48-125=0\), and
\(-10+97-16+25=96\), giving the required right-hand side after division.
\end{solution}

\begin{problem}[Area from side vectors]\label{ex:07:area}
Find the areas of the parallelograms spanned by
(a) \((2,1),(-4,5)\), and (b) \((3,4),(-2,-3)\).
\end{problem}
\begin{solution}
The determinants are \(2(5)-(-4)(1)=14\) and
\(3(-3)-(-2)(4)=-1\). The areas are \(14\) and \(1\); the sign in the
second calculation records orientation, not negative area.
\end{solution}

\begin{problem}[Area from vertices]\label{ex:07:vertices}
Find the parallelogram areas when three vertices are
(a) \((-3,2),(1,4),(-2,-7)\), and
(b) \((2,5),(-1,4),(1,2)\).
\end{problem}
\begin{solution}
Use the first vertex as origin. The side vectors are
\((4,2),(1,-9)\) in (a), giving area \(|-36-2|=38\).
They are \((-3,-1),(-1,-3)\) in (b), giving \(|9-1|=8\).
If a different one of the three vertices is selected, the two difference
vectors change by exchanges and column additions, so the absolute
determinant is unchanged.
\end{solution}

\begin{problem}[Three-dimensional volume]\label{ex:07:volume}
Find the volumes spanned by the following triples:
\begin{enumerate}
\item \((1,1,3),(1,2,-1),(1,4,1)\);
\item \((1,-1,4),(1,1,0),(-1,2,5)\);
\item \((-1,2,1),(2,0,1),(1,3,0)\);
\item \((-2,2,1),(0,1,0),(-4,3,2)\).
\end{enumerate}
\end{problem}
\begin{solution}
Put the vectors in columns and take absolute determinants. The signed values, obtained by expansion in row \(1\), are
for (a)
\(\det\begin{pmatrix}1&1&1\\1&2&4\\3&-1&1\end{pmatrix}
=6-(-11)+(-7)=10\);
for (b),
\[
 \det\begin{pmatrix}1&1&-1\\-1&1&2\\4&0&5\end{pmatrix}
 =5-(-13)+4=22;
\]
for (c),
\[
 \det\begin{pmatrix}-1&2&1\\2&0&3\\1&1&0\end{pmatrix}
 =3+6+2=11.
\]
For (d), subtract twice column \(1\) from column \(3\).
The new third column is \((0,-1,0)^\top\), the negative of column \(2\),
so the determinant is zero. The volumes are \(10,22,11,0\).
\end{solution}

\begin{problem}[Differentiate a determinant]\label{ex:07:derivative}
Let \(B(t),C(t)\in\R^2\) have differentiable entries.
\begin{enumerate}
\item Prove \(\frac{d}{dt}\det(B,C)=\det(B',C)+\det(B,C')\).
\item For twice differentiable \(f,g\), show that
\[
 \frac{d}{dt}\begin{vmatrix}f&g\\f'&g'\end{vmatrix}
 =\begin{vmatrix}f&g\\f''&g''\end{vmatrix}.
\]
\end{enumerate}
\end{problem}
\begin{solution}
Write \(\det(B,C)=b_1c_2-b_2c_1\). The product rule gives
\(b'_1c_2-b'_2c_1+b_1c'_2-b_2c'_1\), the stated two determinants.
For (b), differentiating \(fg'-f'g\) gives
\(f'g'+fg''-f''g-f'g'=fg''-f''g\).
The terms from differentiating the top row cancel because that
determinant has two identical rows.
\end{solution}

\subsection*{Block determinants}
Unless a problem specifies otherwise, write
\(Z=\begin{pmatrix}A&B\\C&D\end{pmatrix}\), where \(A\) is \(m\times m\),
\(D\) is \(n\times n\), \(B\) is \(m\times n\), and \(C\) is \(n\times m\).

\begin{problem}[*Two zero diagonal blocks]\label{ex:07:block-swap}
\begin{enumerate}
\item Show that the block permutation
\(\begin{pmatrix}0&I_m\\I_n&0\end{pmatrix}\) has determinant
\((-1)^{mn}\); its row sizes are \(m,n\) and column sizes \(n,m\).
\item If \(B,C\) are square of orders \(m,n\), show that
\(\det\begin{pmatrix}0&B\\C&0\end{pmatrix}
=(-1)^{mn}\det B\det C\).
\item Explain why the sign is \((-1)^m\) when \(m=n\).
\item If \(B\) is \(m_1\times n_2\), \(C\) is \(m_2\times n_1\),
\(m_1+m_2=n_1+n_2\), and these blocks are not square, determine the
determinant of the same square block matrix.
\end{enumerate}
\end{problem}
\begin{solution}
Moving each of the last \(m\) columns past the first \(n\) columns
takes \(mn\) exchanges and produces the identity, proving (a).
The matrix in (b) is
\(\operatorname{diag}(B,C)\begin{pmatrix}0&I_m\\I_n&0\end{pmatrix}\);
take determinants. For (c), \(m^2-m=m(m-1)\) is even.
In (d), the equality of total sizes and nonsquareness imply either
\(n_2>m_1\) or \(n_1>m_2\).
In the first case \(Bx=0\) has a nonzero solution and \((0,x)\) is in
the kernel of the block matrix; in the second use a nonzero vector in
\(\ker C\). Thus the determinant is zero.
\end{solution}

\begin{problem}[Transpose or inverse in the other block]\label{ex:07:block-transpose}
Find \(\det\begin{pmatrix}0&B\\B^\top&0\end{pmatrix}\) for real
\(B\) of size \(m\times n\), and
\(\det\begin{pmatrix}0&B\\B^{-1}&0\end{pmatrix}\) for invertible
\(m\times m\) \(B\).
\end{problem}
\begin{solution}
If \(m=n\), the first determinant is
\((-1)^m\det B\det B^\top=(-1)^m(\det B)^2\).
If \(m\ne n\), the rectangular case of the preceding problem makes it
zero. The second determinant is
\((-1)^m\det B\det B^{-1}=(-1)^m\).
\end{solution}

\begin{problem}[Remove a Gram block]\label{ex:07:block-gram}
For \(B\) of size \(m\times n\) and square \(D\) of order \(n\), prove
\[
 \det\begin{pmatrix}I_m&B\\B^\top&B^\top B+D\end{pmatrix}=\det D.
\]
\end{problem}
\begin{solution}
Subtract \(B^\top\) times the first block row from the second.
The multiplier \(\begin{pmatrix}I&0\\-B^\top&I\end{pmatrix}\) has
determinant \(1\), and the resulting matrix is
\(\begin{pmatrix}I&B\\0&D\end{pmatrix}\), with determinant \(\det D\).
No invertibility of \(D\) is required.
\end{solution}

\begin{problem}[Another proof of the product theorem]\label{ex:07:block-product}
Use the block triangular formula and elementary row and column
operations to prove \(\det(AB)=\det A\det B\), without using
multiplicativity in the proof.
\end{problem}
\begin{solution}
Starting with \(Q=\begin{pmatrix}A&0\\-I&B\end{pmatrix}\), add the
linear combinations \(A\) of the bottom \(n\) rows to the top \(n\)
rows. These are individual row additions and give
\(R=\begin{pmatrix}0&AB\\-I&B\end{pmatrix}\).
Thus \(\det R=\det Q=\det A\det B\), using the block triangular formula.
Interchange the two column blocks in \(R\), requiring \(n^2\) swaps:
\[
 \det R=(-1)^{n^2}
 \det\begin{pmatrix}AB&0\\B&-I\end{pmatrix}
 =(-1)^{n^2+n}\det(AB)=\det(AB).
\]
The block triangular formula itself follows directly from the permutation
sum, so no step presupposes the desired product theorem.
\end{solution}

\begin{problem}[One zero diagonal block]\label{ex:07:block-zero}
Prove the following formulas under the stated assumptions:
\begin{enumerate}
\item If \(A\) is invertible and \(b,c\in\R^m\),
\(\det\begin{pmatrix}A&b\\c^\top&0\end{pmatrix}
=-\det A\,c^\top A^{-1}b\).
\item If \(A\) is invertible,
\(\det\begin{pmatrix}A&B\\C&0_n\end{pmatrix}
=(-1)^n\det A\,\det(CA^{-1}B)\).
\item If \(D\) is invertible,
\(\det\begin{pmatrix}0_m&B\\C&D\end{pmatrix}
=(-1)^m\det D\,\det(BD^{-1}C)\).
\end{enumerate}
\end{problem}
\begin{solution}
In (b), block elimination leaves diagonal blocks \(A,-CA^{-1}B\);
the scalar case \(n=1\) is (a). In (c), eliminate the upper-right block
using the lower block row, leaving diagonal blocks \(-BD^{-1}C,D\).
A minus sign in an \(r\times r\) matrix contributes \((-1)^r\).
\end{solution}

\begin{problem}[A scalar diagonal block]\label{ex:07:block-scalar}
For \(Z_1=\begin{pmatrix}A&b\\c^\top&\delta\end{pmatrix}\) and
\(Z_2=\begin{pmatrix}\alpha&b^\top\\c&D\end{pmatrix}\), prove
\begin{enumerate}
\item \(\det Z_1=\det A(\delta-c^\top A^{-1}b)\) if \(A\) is invertible;
\item \(\det Z_1=\delta\det(A-\delta^{-1}bc^\top)\) if \(\delta\ne0\);
\item \(\det Z_2=\det D(\alpha-b^\top D^{-1}c)\) if \(D\) is invertible;
\item \(\det Z_2=\alpha\det(D-\alpha^{-1}cb^\top)\) if \(\alpha\ne0\).
\end{enumerate}
\end{problem}
\begin{solution}
For (a), eliminate \(c^\top\) using the first block row; for (b),
eliminate \(b\) using the last row. The respective triangular diagonal
blocks are \(A,\delta-c^\top A^{-1}b\) and
\(A-\delta^{-1}bc^\top,\delta\).
For (c), eliminate \(b^\top\) using the second block row; for (d),
eliminate \(c\) using the first row. The diagonal blocks are
\(\alpha-b^\top D^{-1}c,D\) and \(\alpha,D-\alpha^{-1}cb^\top\).
All elimination multipliers have determinant \(1\).
\end{solution}

\begin{problem}[The other Schur complement and singular blocks]\label{ex:07:schur}
\begin{enumerate}
\item If \(D\) is invertible, prove
\(\det Z=\det D\,\det(A-BD^{-1}C)\).
\item Can both \(A,D\) be singular while \(Z\) is invertible?
\end{enumerate}
\end{problem}
\begin{solution}
Multiply \(Z\) on the left by
\(\begin{pmatrix}I&-BD^{-1}\\0&I\end{pmatrix}\), of determinant \(1\).
The result is \(\begin{pmatrix}A-BD^{-1}C&0\\C&D\end{pmatrix}\),
proving (a). In (b), take scalar blocks \(A=D=0,B=C=1\);
the full matrix has determinant \(-1\) and is its own inverse.
Thus Schur complement formulas must retain the stated invertibility
hypothesis on the block being inverted.
\end{solution}

\begin{problem}[The determinant of a joint Gram matrix]\label{ex:07:joint-gram}
Let \(A\in\R^{N\times k}\), \(B\in\R^{N\times m}\). Prove
\[
 \det\begin{pmatrix}A^\top A&A^\top B\\B^\top A&B^\top B\end{pmatrix}
 =\det(A^\top A)\,
 \det\!\left[B^\top\{I-A(A^\top A)^{-1}A^\top\}B\right]
\]
when \(A\) has independent columns, and prove the corresponding formula
with \(A,B\) interchanged when \(B\) has independent columns.
\end{problem}
\begin{solution}
If \(A^\top Ax=0\), then \(\norm{Ax}^2=x^\top A^\top Ax=0\).
Column independence gives \(x=0\), so \(A^\top A\) is invertible.
Apply the first Schur complement formula to the displayed Gram matrix:
its remaining block is
\(B^\top B-B^\top A(A^\top A)^{-1}A^\top B\), giving the result.
When \(B\) has independent columns the second Schur complement formula
instead gives
\[
 \det(B^\top B)\,
 \det\!\left[A^\top\{I-B(B^\top B)^{-1}B^\top\}A\right].
\]
Only the Gram matrix of the block whose inverse appears must be invertible.
\end{solution}

\begin{problem}[Row-block operations]\label{ex:07:block-operations}
Prove
\[
 \det\begin{pmatrix}EA&EB\\C&D\end{pmatrix}=\det E\det Z
 \quad(E\in\R^{m\times m}),
\]
and
\[
 \det\begin{pmatrix}A&B\\C+EA&D+EB\end{pmatrix}=\det Z
 \quad(E\in\R^{n\times m}).
\]
\end{problem}
\begin{solution}
The two matrices are respectively
\(\operatorname{diag}(E,I_n)Z\) and
\(\begin{pmatrix}I_m&0\\E&I_n\end{pmatrix}Z\).
The multipliers have determinants \(\det E\) and \(1\).
Taking determinants proves both claims, including singular \(E\) in
the first.
\end{solution}

\begin{problem}[Diagonal blocks of an inverse]\label{ex:07:inverse-blocks}
Suppose \(Z\) is invertible and
\(Z^{-1}=\begin{pmatrix}P&Q\\R&S\end{pmatrix}\), with the same block
sizes. Prove
\[
 \det S=\frac{\det A}{\det Z},\qquad
 \det P=\frac{\det D}{\det Z},
\]
even if \(A\) or \(D\) is singular.
\end{problem}
\begin{solution}
The identities in \(ZZ^{-1}=I\) give
\[
 Z\begin{pmatrix}I&Q\\0&S\end{pmatrix}
   =\begin{pmatrix}A&0\\C&I\end{pmatrix},\qquad
 Z\begin{pmatrix}P&0\\R&I\end{pmatrix}
   =\begin{pmatrix}I&B\\0&D\end{pmatrix}.
\]
Taking determinants yields \(\det Z\det S=\det A\) and
\(\det Z\det P=\det D\). Divide by the nonzero \(\det Z\).
No inverse of an individual block has been used.
\end{solution}

\begin{problem}[Two rectangular Gram determinants]\label{ex:07:rectangular-gram}
For \(B\in\R^{m\times n}\), prove
\[
 \det\begin{pmatrix}I_m&B\\B^\top&I_n\end{pmatrix}
 =\det(I_n-B^\top B)=\det(I_m-BB^\top).
\]
\end{problem}
\begin{solution}
Eliminating \(B^\top\) with the upper identity block leaves \(I_n-B^\top B\).
Eliminating \(B\) with the lower identity block leaves \(I_m-BB^\top\).
Both elimination multipliers have determinant \(1\); the other diagonal
block is an identity. This proves both equalities for all rectangular \(B\).
\end{solution}

\begin{problem}[*Two commuting blocks]\label{ex:07:commuting-blocks}
Suppose \(A,B,C,D\) all have the same square order and \(AC=CA\).
Prove
\[
 \det\begin{pmatrix}A&B\\C&D\end{pmatrix}=\det(AD-CB).
\]
\end{problem}
\begin{solution}
First suppose \(A\) is invertible. Commutation implies
\(CA^{-1}=A^{-1}C\). Thus the Schur complement formula gives
\[
 \det A\,\det(D-CA^{-1}B)
 =\det A\,\det(D-A^{-1}CB)=\det(AD-CB).
\]
For a singular \(A\), replace it by \(A+tI\), which still commutes with
\(C\) and is invertible outside finitely many \(t\).
Both sides of the desired equality are polynomials in \(t\) and agree
outside that finite set, so they agree identically. Set \(t=0\).
No commutation involving \(B\) or \(D\) was assumed.
\end{solution}

\begin{problem}[One identity block]\label{ex:07:identity-block}
Prove
\(\det\begin{pmatrix}I_m&B\\C&D\end{pmatrix}=\det(D-CB)\),
including when \(B,C\) are rectangular.
\end{problem}
\begin{solution}
Subtract \(C\) times the first block row from the second.
The determinant is unchanged and the result is
\(\begin{pmatrix}I_m&B\\0&D-CB\end{pmatrix}\), whose determinant is
\(\det(D-CB)\).
\end{solution}

\begin{problem}[A rectangular identity and a determinant update]\label{ex:07:update}
\begin{enumerate}
\item Prove
\[
 \det\begin{pmatrix}I_m&B\\C&I_n\end{pmatrix}
 =\det\begin{pmatrix}I_n&C\\B&I_m\end{pmatrix}.
\]
\item Deduce \(\det(I_m-BC)=\det(I_n-CB)\).
\item If \(A,D\) are invertible of orders \(m,n\), prove
\[
 \det(A+BDC)=\det A\,\det D\,\det(D^{-1}+CA^{-1}B).
\]
\end{enumerate}
\end{problem}
\begin{solution}
Exchanging the two row blocks and then the two column blocks gives the
matrix in (a); the signs multiply to \((-1)^{2mn}=1\).
Eliminating the lower-left or upper-right block gives the two determinants
in (b). For (c), use (b) with \(B\) replaced by \(-A^{-1}B\) and
\(C\) replaced by \(DC\):
\[
 \det(A+BDC)=\det A\,\det(I_m+A^{-1}BDC)
 =\det A\,\det(I_n+DCA^{-1}B).
\]
Factor \(D\) from the final matrix to obtain the result.
\end{solution}

\begin{problem}[A block difference of squares]\label{ex:07:difference}
For square \(A,B\) of the same order, prove
\[
 \det\begin{pmatrix}A&B\\B&A\end{pmatrix}
 =\det(A+B)\det(A-B).
\]
\end{problem}
\begin{solution}
Let \(P=\begin{pmatrix}I&I\\I&-I\end{pmatrix}\); then \(P^{-1}=P/2\).
Direct block multiplication gives
\[
 P^{-1}\begin{pmatrix}A&B\\B&A\end{pmatrix}P
 =\begin{pmatrix}A+B&0\\0&A-B\end{pmatrix}.
\]
Taking determinants cancels \(\det P^{-1}\det P=1\).
The proof requires no assumption \(AB=BA\).
\end{solution}

\begin{problem}[Three diagonal blocks]\label{ex:07:three-blocks}
Let \(A,D,E\) be symmetric square matrices and \(D,E\) invertible.
For conformable \(B,C\), prove
\[
 \det\begin{pmatrix}A&B&C\\B^\top&D&0\\C^\top&0&E\end{pmatrix}
 =\det D\det E\,\det(A-BD^{-1}B^\top-CE^{-1}C^\top).
\]
\end{problem}
\begin{solution}
Group the last two block rows and columns together. The lower-right
block is \(F=\operatorname{diag}(D,E)\), with inverse
\(\operatorname{diag}(D^{-1},E^{-1})\).
Its Schur complement in the first block is
\(A-(B\ C)F^{-1}(B\ C)^\top
=A-BD^{-1}B^\top-CE^{-1}C^\top\).
Apply the second Schur complement formula and
\(\det F=\det D\det E\).
\end{solution}

\begin{problem}[A bordered matrix]\label{ex:07:border}
Let \(A\) be invertible of order \(m\), \(a\in\R^m\), and \(\alpha\ne0\).
Find
\(\det\begin{pmatrix}0_{m\times1}&A\\\alpha&a^\top\end{pmatrix}\).
\end{problem}
\begin{solution}
Column \(1\) has only the entry \(\alpha\), in row \(m+1\).
Cofactor expansion gives
\((-1)^{(m+1)+1}\alpha\det A=(-1)^m\alpha\det A\).
The vector \(a\) has no effect on this determinant.
\end{solution}

